\subsubsection{Select Worked Problems}

\begin{problem}
I roll two dice. Given that the sum is 8, what's the probability that one of the dice shows a 3?
\end{problem}

\begin{solution}
The possible ways to roll a sum of 8 are: $(2, 6), (3, 5), (4, 4), (5, 3), (6, 2)$. Of these, two options have a 3 in them, so the answer is $\frac{2}{5} = 0.4$.
\end{solution}

\begin{problem}
A bag has 3 red and 2 blue balls. You draw one, don't replace it, then draw another. What's P(2nd is red | 1st was blue)?
\end{problem}

\begin{solution}
Once we drew one blue on, the remaining bag composition is 3 red balls and 1 blue one. Then, there is a $\frac{3}{4}$ probability that the second will be red.
\end{solution}


\begin{problem}
60\% of emails are spam. A filter catches 95\% of spam but also flags 3\% of legitimate emails. An email is flagged. What's the probability it's actually spam?
\end{problem}

\begin{solution}
We want $P(Spam | Flagged)$. The likelihood $P(Flagged | Spam)$ is given as $0.95$, the prior $P(Spam)$ is $0.6$, and the marginal probability $P(Flagged)$ is $0.95 \cdot 0.6 + 0.03 \cdot 0.4 = 0.582$.

Then the posterior is
\[
P(Spam | Flagged) = \frac{P(Flagged | Spam) \times P(Spam)}{P(Flagged)} = \frac{0.95 \times 0.6}{0.582} \approx 0.979
\]


\end{solution}


\begin{problem}
You have 2 coins. Coin A is fair (50\% heads). Coin B is biased (75\% heads). You pick a coin at random and flip it. It lands heads. What's the probability you picked coin B?
\end{problem}

\begin{solution}
We are looking for the posterior, $P(B|H)$.

The priors are $P(A) = P(B) = 0.5$. The likelihoods are $P(H|A) = 0.5$ and $P(H|B) = 0.75$. Then $P(H) = P(H|A) \cdot P(A) + P(H|B) \cdot P(B) = 0.5 \cdot 0.5 + 0.75 \cdot 0.5 = 0.625$. Then finally
\[
P(B|H) = \frac{P(H|B) \cdot P(B)}{P(H)} = \frac{0.75 \cdot 0.5}{0.625} = 0.6
\]
After flipping and landing on heads, our belief that we picked coin B increased from $0.5$ to $0.6$.
\end{solution}

\begin{problem}
Same as 1.31. The chosen coin is flipped twice and get heads both times. Now what's $P(B)$?
\end{problem}

\begin{solution}
The only difference is our priors are $P(A) = 0.4$ and $P(B) = 0.6$. Then we update $P(H) = P(H|A) \cdot P(A) + P(H|B) \cdot P(B) = 0.5 \cdot 0.4 + 0.75 \cdot 0.6 = 0.65$. Finally:
\[
P(B|H) = \frac{P(H|B) \cdot P(B)}{P(H)} = \frac{0.75 \cdot 0.6}{0.65} = \frac{9}{13} \approx 0.692
\]
\end{solution}

\newpage

\begin{problem}
A family has 2 children. Given that at least one child is a boy, what's the probability both are boys?
\end{problem}

\begin{solution}
We can answer this by considering the restriction over the sample space caused by the new information. Originally, for two children there would be four possibilities, of which 1 has two boys. After we know that at least one child is a boy, only three options remain $(B,G), (B,B), (G, B)$. Thus, the probability both are boys is $\frac{1}{3}$.
\end{solution}

\begin{problem}
Three cards: one is red on both sides, one is blue on both sides, one is red on one side and blue on the other. You pick a card at random and see a red side. What's the probability the other side is also red?
\end{problem}

\begin{solution}
The natural first instinct is to treat the cards as the sample space, as with the births problem: restrict to cards consistent with seeing a red face, then count outcomes among those. This fails here because a card is not a single unit of observation. When a card has two sides of the same color, ``seeing red'' can happen in more than one way for that card. There's no such internal multiplicity in the births problem.

So take the six sides, not the three cards, as the equally likely outcomes. Three of these six are red: two belong to the $RR$ card, one belongs to the $RB$ card. Since each red side was equally likely to be the one we observed, and two of the three red sides have a red back:
\[
P(\text{other side red} \mid \text{observed red}) = \frac{2}{3}.
\]
\end{solution}


\begin{problem}
You have a 6-sided die. You don't know if it's fair or loaded. A fair die shows each number with probability 1/6. The loaded die shows 6 with probability 1/2 and each other number with probability 1/10. You pick a die at random (50-50) and roll a 6. What's P(loaded)?
\end{problem}

\begin{solution}
By Bayes' theorem, with $P(\text{loaded}) = P(\text{fair}) = 0.5$:
\[
P(\text{loaded} \mid 6) = \frac{P(6 \mid \text{loaded}) \, P(\text{loaded})}{P(6)} = \frac{0.5 \times 0.5}{0.5 \times 0.5 + \tfrac{1}{6} \times 0.5} = \frac{0.25}{1/3} = 0.75
\]

Alternatively, work in odds instead of probabilities. Dividing the Bayes' theorem expressions for $P(\text{loaded} \mid 6)$ and $P(\text{fair} \mid 6)$ cancels $P(6)$ from both, leaving:
\[
\underbrace{\frac{P(\text{loaded} \mid 6)}{P(\text{fair} \mid 6)}}_{\text{posterior odds}} = \underbrace{\frac{P(6 \mid \text{loaded})}{P(6 \mid \text{fair})}}_{\text{likelihood ratio}} \times \underbrace{\frac{P(\text{loaded})}{P(\text{fair})}}_{\text{prior odds}}
\]
So the posterior odds are just the prior odds scaled by the likelihood ratio, no need to compute $P(6)$ at all. Here, the prior odds are $\frac{P(\text{loaded})}{P(\text{fair})} = \frac{0.5}{0.5} = 1$, and the likelihood ratio is $\frac{P(6 \mid \text{loaded})}{P(6 \mid \text{fair})} = \frac{0.5}{1/6} = 3$. So the posterior odds are $1 \times 3 = 3$, i.e. $3:1$ in favor of loaded, giving $P(\text{loaded} \mid 6) = 3/4 = 0.75$.

In short, a six is three times more likely to appear if the die is loaded, so each 6 observed multiplies the odds by 3. If we observed \textit{another} 6, the odds would update again: $3:1 \rightarrow 9:1$, giving $P(\text{loaded}) = \frac{9}{10}$.

\end{solution}
